\section{Results}
All results stem from the datasets generated during the practical-application
 of Bayesian optimisation in the domain of thermoset composite property
 development.
Practical experimental results were split into those derived from the
 single-objective experiments and the singular multi-objective optimisation.

\subsection{Single-Objective Experiments}
\label{subsec:UniObjRes}
Two single-objective Bayesian optimisation experiments were performed, each
 using a different objective variable for sequential sampling: Young's 
 modulus \(E\) or yield strength \(\sigma_{\mathrm{y}}\).
Initial Gaussian process priors developed for both experiments were trained
 on the same eight quasirandomly sampled (Sobol) points \(\mathbf{X}\) (from
 across the parameter space \(\mathcal{X}\)) and their evaluations
 \(\mathbf{Y}\).
Both experiments proceeded sequentially with a single sample evaluated at 
 each Bayesian iteration.
The sample limit was set to 24 (including the
 prior-forming dataset), and the optimisation was one of objective metric
 maximisation.

The objective-naive MIPT sampling technique was used to sample in a purely 
 exploratory fashion, which provided a benchmark against which Bayesian 
 optimisation routine progression could be gauged.
Given its objective-naivety, purely exploratory nature, and the fact that
 it was initiated on the same prior-forming dataset as the Bayesian routines,
 the MIPT routine was only run once (samples having both Young's
 modulus and yield strength evaluated).
The resulting dataset could be directly compared with both the single- 
 and multi-objective Bayesian sequential routines (see Section~\ref{subsec:MultiObjRes}).

\paragraph{Young's Modulus Experiment}
Figure~\ref{fig:OptProgYM} compares the single-objective optimisation 
 progression of both Bayesian and exploratory sequential routines in the 
 maximisation of Young's modulus.
The one-shot quasirandomly-sampled initial dataset provided a baseline
 maximum of 4274~\unit{\mega\pascal} and a mean of 3052~\unit{\mega\pascal}.
On reaching the sample limit, the single-objective Bayesian technique had 
 exceeded both these baselines (max 5012~\unit{\mega\pascal}, mean 
 3396~\unit{\mega\pascal}) whilst the exploratory sampler only exceeded the 
 maximum baseline (max 4851~\unit{\mega\pascal}, mean 3046~\unit{\mega\pascal}).
Although the single-objective Bayesian routine ultimately discovered a sample exhibiting the
 highest Young's modulus within this experiment, it is interesting that the
 exploratory routine discovered a greater number of points surpassing the
 baseline maximum set by the initial dataset.
The exploratory sampling routine was the first to locate a sample improving
 upon the baseline maximum: sample 15 with a Young's modulus of 
 4851~\unit{\mega\pascal}.
Subsequently, at sample 17, the Bayesian single-objective routine
 discovered the experiment's highest Young's modulus of
 5012~\unit{\mega\pascal}.

\noindent
\begin{minipage}{0.485\textwidth}
    \includegraphics[width=\textwidth]{SOBO-YM-Progression.png}
    \captionsetup{format=custom,width=\textwidth,skip=0pt}
    \captionof{figure}{
        \footnotesize{
            Progression of a the sequential optimisation experiment concerning
            the maximisation of Young's modulus.
            Initial one-shot samples drawn quasirandomly are presented
             separately from the two sequential techniques which they prime.
            These two sequential techiques are either objective-naive (MIPT)
             or single-objective (Bayesian).
        }
    }
    \label{fig:OptProgYM}
\end{minipage}

\paragraph{Yield Point Experiment}
Figure~\ref{fig:OptProgYS} compares optimisation progression of both the
 Bayesian single-objective and exploratory sequential routines with respect 
 to the maximisation of the yield point.
The one-shot quasirandomly-sampled initial dataset provided a baseline
 maximum yield point of 137~\unit{\mega\pascal} and a mean of 123~\unit{\mega\pascal}.
On reaching the sample limit, the Bayesian technique had surpassed the
 baseline maximum (143~\unit{\mega\pascal}) and remained level with respect
 to the mean (123~\unit{\mega\pascal}), whilst the exploratory routine
 had surpassed the baseline maximum (140~\unit{\mega\pascal}) and failed to
 improve upon the baseline mean (118~\unit{\mega\pascal}).
In this experiment both the Bayesian single-objective and exploratory routines discovered the
 same number of points surpassing the baseline maximum, although the
 Bayesian technique ultimately uncovered the sample with the highest yield
 strength.
The exploratory sampling technique was the first to discover a point
 improving upon the maximum baseline; with this happening on sample 15,
 with a yield strength of 140~\unit{\mega\pascal}.
Subsequently, the Bayesian sequential technique improves upon the maximum
 baseline on sample 18 (140~\unit{\mega\pascal})
 and surpasses the exploratory routine on sample 21, with a final yield 
 strength of 143~\unit{\mega\pascal}.

\noindent
\begin{minipage}{0.485\textwidth}
    \includegraphics[width=\textwidth]{SOBO-YS-Progression.png}
    \captionsetup{format=custom,width=\textwidth,skip=0pt}
    \captionof{figure}{
        \footnotesize{
            Progression of a sequential optimisation experiment concerning
            the maximisation of yield strength.
            Initial one-shot samples drawn quasirandomly are presented
             separately from the two sequential techniques which they prime.
            These two sequential techiques are either objective-naive (MIPT)
             or single-objective (Bayesian).
        }
    }
    \label{fig:OptProgYS}
\end{minipage}

\paragraph{Comparison of Bayesian and MIPT Techniques}
In both single-objective optimisation experiments, the Bayesian sequential
 routines outperformed the exploratory sampling routine in terms of
 discovery of samples approaching an experimental global maximum within the 
 alotted sample limit.
Bayesian techniques benefit from adaptive sampling based on the
 balancing of both exploitative and exploratory behaviour, the former 
 exhibited clearly in the non-decreasing trends shown by the Bayesian 
 running means seen in both Figures~\ref{fig:OptProgYM}\&\ref{fig:OptProgYS}.
These results support the hypothesis that the Bayesian sampling
 routine was more effective in maximising the objective property during the
 course of a sequential optimisation, and although the margin of improvement
 was slim they remain significant.
Further support may be garnered from repeat experiments of either a physical 
 or computational nature.


However, in both experiments, the exploratory sampling routine 
 outperformed the Bayesian single-objective technique in terms of sample 
 efficiency: the exploratory technique improved upon the baseline maxima 
 (\(E=4274\)~\unit{\mega\pascal} and \(\sigma_{y}=137\)~\unit{\mega\pascal}) 
 of both experiments on drawing sample 15 for evaluation (\(E=4851\) and
 \(\sigma_{y}=140\)).
The fact that this occurred on the 15\textsuperscript{th} sample for both
 experiments is trivially explained by the fact that this was the same 
 exploratory routine being performed.
As opposed to the two single-objective routines seen in Figure~\ref{fig:OptProgYM}
 and Figure~\ref{fig:OptProgYS} which were sequentially drawing unique sample in attempts
 to maximise different properties, the objective-naive MIPT sampler was 
 only run once with each sequential sample characterised in terms of both
 target variables.
This is why the especially well-performing sample 15 appears in both 
 progression plots (Figs.~\ref{fig:OptProgYM}/\ref{fig:OptProgYS}).
A hypothesis that the exploratory sampler more efficiently improves upon
 the baseline maximum than the Bayesian routine would be weakly supported
 by the above findings given that both circumstances were precipitated from
 the same data point, but further repeat studies in the physical or
 computational domain would aid in the supporting or negating this
 conjecture.

\paragraph{Styrenated Young's Modulus Experiment}
An experiment identical to that of the single-objective Bayesian optimisation
 for Young's modulus was carried out with the conventional reactive diluent
 styrene in place of the bio-based dimethyl itaconate.
Figure~\ref{fig:OptProgYSwSTY} shows an optimisation progression plot for
 the styrene-based system.
The single-objective Bayesian optimisation technique shows strong
 performance in terms of both sample efficiency and scale of improvement
 above the baseline.
On conclusion of this sequential optimisation, the putative global
 optimum (sample 18) exhibited a compressive Young's modulus of
 7,271~MPa.
Interestingly, the analogous bio-based experiment produced a sample
 (sample 17 in Fig.~\ref{fig:OptProgYM}) exhibiting 69\% the Young's
 modulus of this styrenated putative global optimum.
Furthermore, when the relative crosslink densities between diitaconate
 and styrene based resins of 73\unit{\percent} reported by Panic et al.
 are considered \autocite{Panic2017}, it becomes apparent that this
 short bout of optimisation has resulted in samples approaching the
 currently known bounds of mechanical performance.

\noindent
\begin{minipage}{0.485\textwidth}
    \includegraphics[width=\textwidth]{SOBO-YM-Progression-STY.png}
    \captionsetup{format=custom,width=\textwidth,skip=0pt}
    \captionof{figure}{
        \footnotesize{
            Progression of a sequential optimisation experiment concerning
            the maximisation of Young's modulus with a styrene-based
            thermoset.
        }
    }
    \label{fig:OptProgYSwSTY}
\end{minipage}

\subsection{Multi-Objective Experiment}
\label{subsec:MultiObjRes}
In the multi-objective setting, a Bayesian sequential routine was performed
 which made use of the noisy expected hypervolume improvement acquisition
 function to sample-efficiently identify the Pareto front within an objective
 space comprised of both Young's modulus and yield strength.
This dual-objective optimisation was facilitated by two separate Gaussian
 processes modelling the two different objective variables.
Both were trained on 56 samples concatenated from the re-analysed raw data 
 of the single-objective experiments described in Section~\ref{subsec:UniObjRes}.
The order of concatenation (and re-indexing of sample numbers) was as
 follows: one-shot objective-naive quasirandom sampler (samples 1-8),
 sequential exploratory objective-naive routine (9-24), sequential Bayesian
 yield strength maximising routine (25-40), and the sequential Bayesian
 Young's modulus maximising routine (41-56).
This prior-forming dataset \(\mathcal{D}\) yielded an initial Pareto front
 \(\mathcal{P}\) comprised of three samples: two having been discovered
 as maxima by Bayesian sequential routines in their two respective
 single-objective experiments with yield strength (sample 37) an
 Young's modulus (sample 49) as their objectives.
Whilst the third (sample 15) was derived from the exploratory technique as a
 point which performed well in both objectives of interest.
With an ultimate sample limit of 72 and a single sample drawn per Bayesian
 iteration, the sequential multi-objective Bayesian routine proceeded to
 develop the Pareto front over 15 further samples (57-71).

All samples both prior-forming and sequentially-obtained are depicted
 in Figure~\ref{fig:OptMO}, with the initial Pareto front marked as a
 dashed black line (Fig.\ref{fig:OptMO}A).
Given the initial Pareto front and considering only the direct evaluations
 obtained for samples it can be said that five samples were discovered
 which improved upon the baseline described by the Pareto front
 (samples 62,66,68,69,71).
Under progression of the multi-objective routine, the first novel pareto
 optimal point which improved the Pareto front was the 62\textsuperscript{nd}
 sample which dominated the 37\textsuperscript{th} sample shifting the
 boundaries of the developing Pareto front (Fig.\ref{fig:OptMO}B).
Subsequently, the sample 66 evaluated
 (\(E=5043\)~\unit{\mega\pascal} and \(\sigma_{y}=153\)~\unit{\mega\pascal})
 presented a marginal improvement upon both the Young's modulus of
 sample 49 (\(E=5012\)~\unit{\mega\pascal}) and the yield strength of sample
 62 (\(\sigma_{y}=150\)~\unit{\mega\pascal}), collapsing the Pareto front to
 a single Pareto optimal point.
Finally, the sample 68 dominated sample 66
 producing the putative globally Pareto optimal solution 
 (\(E=5458\)~\unit{\mega\pascal} and \(\sigma_{y}=165\)~\unit{\mega\pascal}),
 which may be contrasted starkly with the anti-Pareto front comprised of
 the 14\textsuperscript{th} (\(E=1114\)~\unit{\mega\pascal} and
 \(\sigma_{y}=62\)~\unit{\mega\pascal}) and 25\textsuperscript{th} 
 (\(E=813\)~\unit{\mega\pascal} and \(\sigma_{y}=105\)~\unit{\mega\pascal})
 samples.
How the acquisition function likely navigated the parameter space and came
 to the samples it discovered has been reviewed from the
 perspective of the final Gaussian processes developed and is appended in
 the supporting information.

\noindent
\begin{minipage}{0.485\textwidth}
    \includegraphics[width=\textwidth]{MOBO-Progression.png}
    \captionsetup{format=custom,width=\textwidth,skip=0pt}
    \captionof{figure}{
        \footnotesize{
            Results from a multi-Objective Bayesian optimisation routine.
            Bayesian confidence intervals were a single standard
            deviation and sampled from the final Gaussian processes.
            (A) Initial Pareto front as a dashed black line.
            (B) Pareto front improvement as a dashed green line.
            YS, Yield strength.
            YM, Young's modulus.
            SOBO, Single-objective Bayesian optimisation.
            MOBO, Multi-objective Bayesian optimisation.
        }
    }
    \label{fig:OptMO}
\end{minipage}

Figure~\ref{fig:MOBOSampleComp} presents the strong relation between
 initiating agent content maximisation and improvements to sample optimality:
 the average content is 13\% for the anti-Pareto front, 44\% for the initial
 Pareto front, and 50\% for the putative global optimum.
Furthermore, the initiating agent's composition also described a somewhat
 weaker relation between dibenzoyl peroxide marginalisation and sample
 optimality: anti-Pareto front 13.3\%, initial Pareto front 12.3\%, and 8\%
 for the putative global optimum sample.
No analagously clear trend existed with respect to the unsaturated polyester
 resin composition.

\noindent
\begin{minipage}{0.485\textwidth}
    \includegraphics[width=\textwidth]{MOBO-SampleCompositions.png}
    \captionsetup{format=custom,width=\textwidth,skip=0pt}
    \captionof{figure}{
        \footnotesize{
            Compositional breakdown for samples of interest including the
            initial Pareto front, the final anti-Pareto front, and the final
            Pareto front.
        }
    }
    \label{fig:MOBOSampleComp}
\end{minipage}